\section{Niveau 1}

\subsection{Les variables : une case mémoire nommée}

\begin{UPSTIexemples}[La variable : une case étiquetée]{}
    Une variable, c'est comme une variable Scratch : une \textbf{case mémoire}, portant une \textbf{étiquette} (son nom) et contenant une \textbf{valeur} qui peut changer au cours du programme. En C, cette case a en plus un \textbf{type} fixé une fois pour toutes, qui indique la nature de ce qu'elle contient.

    \begin{center}
        \begin{tabular}{|c|c|c|}
            \hline
            \textbf{Étiquette (nom)} & \textbf{Type}  & \textbf{Contenu (valeur)} \\
            \hline
            \texttt{nbEleves}        & \texttt{int}   & \texttt{24}               \\
            \hline
        \end{tabular}
    \end{center}

    En C, cette case se crée et se remplit ainsi :
    \begin{lstlisting}[language=C]
int nbEleves;   // declaration : on cree la case, de type int
nbEleves = 24;  // affectation : on range 24 dans la case

// ou, en une seule ligne (declaration + initialisation) :
int nbEleves = 24;
    \end{lstlisting}
\end{UPSTIexemples}

\begin{UPSTIexercice}{Repérer les variables}
    On considère le programme suivant, qui affiche des informations sur un capteur de température.
    \begin{lstlisting}[language=C]
#include <stdio.h>

int main(void)
{
    int numeroCapteur = 12;
    char etat = 'M';
    float temperature;
    temperature = 21.5;
    numeroCapteur = numeroCapteur + 1;

    printf("Capteur %i, etat %c, temperature %.1f\n", numeroCapteur, etat, temperature);
}
    \end{lstlisting}
    \UPSTIquestion{\textbf{Entourer} toutes les lignes où une variable est \textbf{déclarée}.}
    \UPSTIquestion{\textbf{Souligner} toutes les lignes où une variable est \textbf{affectée} (reçoit une valeur).}
    \UPSTIquestion{Compléter le tableau ci-dessous avec le type de chaque variable et sa valeur juste avant l'appel à \texttt{printf}.}

    \begin{center}
        \begin{tabular}{|l|l|l|}
            \hline
            \textbf{Nom de la variable} & \textbf{Type} & \textbf{Valeur avant le \texttt{printf}} \\
            \hline
             &  &  \\[0.6cm]
            \hline
             &  &  \\[0.6cm]
            \hline
             &  &  \\[0.6cm]
            \hline
        \end{tabular}
    \end{center}
\end{UPSTIexercice}

\begin{UPSTIprofOnlyEnv}
    \begin{UPSTIcorrectionP}{Repérer les variables}
        Les déclarations sont les lignes \texttt{int numeroCapteur = 12;}, \texttt{char etat = 'M';} et \texttt{float temperature;} (cette dernière est une déclaration sans initialisation).

        Les affectations sont les lignes \texttt{temperature = 21.5;} et \texttt{numeroCapteur = numeroCapteur + 1;} (cette dernière lit l'ancienne valeur de \texttt{numeroCapteur} pour calculer la nouvelle).

        \begin{center}
            \begin{tabular}{|l|l|l|}
                \hline
                \textbf{Nom}      & \textbf{Type}   & \textbf{Valeur avant le \texttt{printf}} \\
                \hline
                numeroCapteur     & int             & 13                                       \\
                \hline
                etat              & char            & 'M'                                      \\
                \hline
                temperature       & float           & 21.5                                     \\
                \hline
            \end{tabular}
        \end{center}
    \end{UPSTIcorrectionP}
\end{UPSTIprofOnlyEnv}


\begin{UPSTIexercice}{Trace d'exécution : la caisse d'un distributeur}
    Un distributeur de boissons exécute les lignes suivantes lorsqu'un client insère des pièces. Il ne rend la monnaie qu'avec des pièces de 10 centimes, et affiche à part ce qu'il ne peut pas rendre avec ce type de pièce.
    \begin{lstlisting}[language=C]
int credit = 0;
int prix = 85;
int monnaie;
int nbPieces10;
int centimesRestants;

credit = credit + 50;            // (1) le client insere 50 centimes
credit = credit + 50;            // (2) le client insere encore 50 centimes
monnaie = credit - prix;         // (3)
nbPieces10 = monnaie / 10;       // (4) nombre de pieces de 10 centimes a rendre
centimesRestants = monnaie % 10; // (5) ce qui ne peut pas etre rendu en pieces de 10
credit = 0;                      // (6) la machine rend la monnaie, le credit repart a 0
    \end{lstlisting}
    \UPSTIquestion{Rappeler ce que calculent, en général, les opérateurs \texttt{/} et \texttt{\%} entre deux entiers (voir les encarts page~\pageref{info:division} et page~\pageref{info:modulo}).}
    \UPSTIquestion{Compléter le tableau de trace ci-dessous : pour chaque ligne numérotée, indiquer la valeur de chaque variable \textbf{juste après} l'exécution de cette ligne.}

    \begin{center}
        \begin{tabular}{|c|c|c|c|c|}
            \hline
            \textbf{Ligne} & \texttt{credit} & \texttt{monnaie} & \texttt{nbPieces10} & \texttt{centimesRestants} \\
            \hline
            avant (1)      & 0               & ?                & ?                   & ?                         \\
            \hline
            (1)            &                 &                  &                     &                           \\[0.5cm]
            \hline
            (2)            &                 &                  &                     &                           \\[0.5cm]
            \hline
            (3)            &                 &                  &                     &                           \\[0.5cm]
            \hline
            (4)            &                 &                  &                     &                           \\[0.5cm]
            \hline
            (5)            &                 &                  &                     &                           \\[0.5cm]
            \hline
            (6)            &                 &                  &                     &                           \\[0.5cm]
            \hline
        \end{tabular}
    \end{center}
    \UPSTIquestion{Combien de pièces de 10 centimes le distributeur rend-il au client, et quel montant, en centimes, ne peut-il pas rendre avec ce type de pièce ?}
    \UPSTIquestion{\texttt{nbPieces10} et \texttt{centimesRestants} changent-ils de valeur après la ligne (6) ? Justifier.}
\end{UPSTIexercice}

\begin{UPSTIprofOnlyEnv}
    \begin{UPSTIcorrectionP}{Trace d'exécution : la caisse d'un distributeur}
        Entre deux entiers, \texttt{/} calcule le quotient de la division euclidienne (la partie décimale est perdue), et \texttt{\%} calcule le reste de cette même division.

        \begin{center}
            \begin{tabular}{|c|c|c|c|c|}
                \hline
                \textbf{Ligne} & \texttt{credit} & \texttt{monnaie}  & \texttt{nbPieces10} & \texttt{centimesRestants} \\
                \hline
                avant (1)      & 0               & (indéfinie)       & (indéfinie)         & (indéfinie)               \\
                \hline
                (1)            & 50              & (indéfinie)       & (indéfinie)         & (indéfinie)               \\
                \hline
                (2)            & 100             & (indéfinie)       & (indéfinie)         & (indéfinie)               \\
                \hline
                (3)            & 100             & 15                & (indéfinie)         & (indéfinie)               \\
                \hline
                (4)            & 100             & 15                & 1                   & (indéfinie)               \\
                \hline
                (5)            & 100             & 15                & 1                   & 5                         \\
                \hline
                (6)            & 0               & 15                & 1                   & 5                         \\
                \hline
            \end{tabular}
        \end{center}
        Le distributeur rend \textbf{1 pièce de 10 centimes}, et il lui reste \textbf{5 centimes} qu'il ne peut pas rendre avec ce type de pièce (\texttt{15 = 1 $\times$ 10 + 5}).

        \texttt{nbPieces10} et \texttt{centimesRestants} ne changent pas après la ligne (6) : cette ligne ne modifie que \texttt{credit}, chaque variable étant une case mémoire indépendante des autres.
    \end{UPSTIcorrectionP}
\end{UPSTIprofOnlyEnv}


\begin{UPSTIexercice}{Choisir le bon type}
    On souhaite programmer un petit boîtier de supervision qui doit stocker les informations suivantes :
    \begin{itemize}
        \item le nombre de capteurs actuellement connectés au boîtier ;
        \item la tension mesurée aux bornes d'une LED (en volts) ;
        \item l'état d'un interrupteur (ouvert ou fermé) ;
        \item la référence d'un composant, du type \texttt{"R220"} ;
        \item la lettre correspondant à la gamme du boîtier (\texttt{'A'}, \texttt{'B'} ou \texttt{'C'}).
    \end{itemize}
    \UPSTIquestion{Pour chaque information, indiquer le type C le plus adapté parmi \texttt{int}, \texttt{float}, \texttt{char}, \texttt{string} et \texttt{bool}, et justifier brièvement.}
    \UPSTIquestion{Écrire, pour chacune, une ligne de déclaration avec une initialisation cohérente.}

\end{UPSTIexercice}

\begin{UPSTIprofOnlyEnv}
    \begin{UPSTIcorrectionP}{Choisir le bon type}
        \begin{lstlisting}[language=C]
// Nombre de capteurs connectes : un compte, toujours entier
int nbCapteurs = 4;

// Tension aux bornes d'une LED : une grandeur mesuree, avec decimales
float tension = 1.8;

// Etat d'un interrupteur : deux possibilites seulement -> booleen
bool interrupteurFerme = true;

// Reference d'un composant : du texte
string reference = "R220";

// Gamme du boitier : un seul caractere
char gamme = 'B';
        \end{lstlisting}
    \end{UPSTIcorrectionP}
\end{UPSTIprofOnlyEnv}


\begin{UPSTIexercice}{Prédire l'affichage}
    On exécute le programme suivant.
    \begin{lstlisting}[language=C]
int nb = 8;
float tension = 3.3;
char gamme = 'B';

printf("Nombre de capteurs : %i\n", nb);
printf("Tension mesuree : %.1f V\n", tension);
printf("Gamme : %c\n", gamme);
printf("Tension precise : %.4f V\n", tension);
    \end{lstlisting}
    \UPSTIquestion{Prédire, ligne par ligne, ce qui s'affiche dans la console.}
    \UPSTIquestion{La dernière ligne affiche \texttt{tension} avec 4 décimales. Le résultat sera-t-il exactement \texttt{3.3000} ? Justifier en une phrase à l'aide du cours sur l'imprécision des flottants.}
\end{UPSTIexercice}

\begin{UPSTIprofOnlyEnv}
    \begin{UPSTIcorrectionP}{Prédire l'affichage}
        Affichage attendu :
        \begin{lstlisting}[language=bash,style=console]
Nombre de capteurs : 8
Tension mesuree : 3.3 V
Gamme : B
Tension precise : 3.3000
        \end{lstlisting}
        En toute rigueur, un \texttt{float} ne peut pas représenter \texttt{3.3} de façon exacte en binaire : la valeur réellement stockée est une approximation très proche de 3.3 (par exemple \texttt{3.29999995...}). Avec seulement 4 décimales affichées, l'arrondi masque cette imprécision et on obtient bien \texttt{3.3000} à l'affichage ; mais avec davantage de décimales (\texttt{\%.10f} par exemple), l'imprécision deviendrait visible, comme dans le cours.
    \end{UPSTIcorrectionP}
\end{UPSTIprofOnlyEnv}


\subsection{Les structures conditionnelles : un aiguillage}

\begin{UPSTIexemples}[Un aiguillage simple]{}
    Une structure conditionnelle fonctionne comme un aiguillage : selon la valeur d'une condition (vraie ou fausse), le programme n'exécute \textbf{qu'un seul} des blocs proposés, jamais plusieurs, jamais aucun (s'il y a un \texttt{SINON}).

    On surveille une variable booléenne \texttt{alarme}. Si elle vaut vrai, on doit afficher \og Evacuation \fg{}, sinon \og RAS \fg.

    \noindent
    \begin{minipage}[t]{0.48\linewidth}
        \textbf{Pseudo-code}
        \begin{verbatim}
SI alarme == vrai ALORS
    AFFICHER "Evacuation"
SINON
    AFFICHER "RAS"
FIN SI
        \end{verbatim}
    \end{minipage}\hfill
    \begin{minipage}[t]{0.48\linewidth}
        \textbf{C}
        \begin{lstlisting}[language=C]
if (alarme)
{
  printf("Evacuation\n");
}
else
{
  printf("RAS\n");
}
        \end{lstlisting}
    \end{minipage}

    \textbf{Trace pour \texttt{alarme = faux} :} la condition \texttt{alarme == vrai} est fausse, donc c'est le bloc \texttt{SINON} qui s'exécute : le programme affiche \og RAS \fg. Le bloc \texttt{ALORS} n'est jamais atteint.
\end{UPSTIexemples}

\begin{UPSTIexercice}{Repérer les blocs}
    On considère le programme suivant, qui classe une note sur 20.
    \begin{lstlisting}[language=C]
int note = 14;

if (note < 8)
{
    printf("Insuffisant\n");
}
else if (note < 12)
{
    printf("Passable\n");
}
else
{
    printf("Satisfaisant\n");
}
    \end{lstlisting}
    \UPSTIquestion{\textbf{Entourer} chacune des conditions testées (il y en a deux).}
    \UPSTIquestion{\textbf{Souligner} d'un trait chaque bloc de code associé à une branche (il y en a trois), en utilisant une couleur différente par bloc si possible.}
    \UPSTIquestion{Ces trois blocs sont-ils \textbf{mutuellement exclusifs}, c'est-à-dire qu'un seul peut s'exécuter pour une valeur de \texttt{note} donnée ? Justifier.}
    \UPSTIquestion{Pour \texttt{note = 14}, numéroter dans l'ordre où elles sont réellement évaluées les étapes suivantes : \emph{test de \texttt{note < 8}}, \emph{test de \texttt{note < 12}}, \emph{affichage de \og Satisfaisant \fg}.}
\end{UPSTIexercice}

\begin{UPSTIprofOnlyEnv}
    \begin{UPSTIcorrectionP}{Repérer les blocs}
        Les deux conditions sont \texttt{note < 8} et \texttt{note < 12}. Les trois blocs sont ceux qui affichent respectivement \og Insuffisant \fg, \og Passable \fg{} et \og Satisfaisant \fg.

        Oui, les trois blocs sont mutuellement exclusifs : la structure \texttt{if / else if / else} garantit qu'un seul bloc s'exécute, car dès qu'une condition est vraie, les suivantes ne sont même plus testées.

        Pour \texttt{note = 14} : (1) test de \texttt{note < 8} $\to$ faux ; (2) test de \texttt{note < 12} $\to$ faux ; (3) aucune condition n'étant vraie, c'est le bloc \texttt{else} qui s'exécute directement, sans nouveau test : affichage de \og Satisfaisant \fg.
    \end{UPSTIcorrectionP}
\end{UPSTIprofOnlyEnv}


\begin{UPSTIexercice}{Feu tricolore simplifié (exercice guidé)}
    On modélise un feu piéton par une variable \texttt{couleur} qui vaut \texttt{'R'} (rouge) ou \texttt{'V'} (vert). On veut afficher \og Attendre \fg{} si le feu est rouge, \og Traverser \fg{} sinon.

    \UPSTIquestion{Compléter les pointillés du pseudo-code ci-dessous.}
    \begin{verbatim}
SI couleur == ......... ALORS
    AFFICHER .....................
SINON
    AFFICHER .....................
FIN SI
    \end{verbatim}
    \UPSTIquestion{Compléter les pointillés du code C ci-dessous, qui doit produire le même résultat.}
    \begin{lstlisting}[language=C]
if (couleur == .........)
{
    printf(".....................\n");
}
else
{
    printf(".....................\n");
}
    \end{lstlisting}
\end{UPSTIexercice}

\begin{UPSTIprofOnlyEnv}
    \begin{UPSTIcorrectionP}{Feu tricolore simplifié (exercice guidé)}
        \begin{verbatim}
SI couleur == 'R' ALORS
    AFFICHER "Attendre"
SINON
    AFFICHER "Traverser"
FIN SI
        \end{verbatim}
        \begin{lstlisting}[language=C]
if (couleur == 'R')
{
    printf("Attendre\n");
}
else
{
    printf("Traverser\n");
}
        \end{lstlisting}
    \end{UPSTIcorrectionP}
\end{UPSTIprofOnlyEnv}


\begin{UPSTIexercice}{Chauffage automatique (exercice autonome)}
    Un radiateur programmable doit afficher \og Chauffage ON \fg{} si la température mesurée \texttt{temperature} est strictement inférieure à \texttt{19.0}, et \og Chauffage OFF \fg{} sinon.
    \UPSTIquestion{Écrire, entièrement seul, le pseudo-code correspondant.}

    \UPSTIquestion{Traduire ce pseudo-code en langage C.}

\end{UPSTIexercice}

\begin{UPSTIprofOnlyEnv}
    \begin{UPSTIcorrectionP}{Chauffage automatique (exercice autonome)}
        \begin{verbatim}
SI temperature < 19.0 ALORS
    AFFICHER "Chauffage ON"
SINON
    AFFICHER "Chauffage OFF"
FIN SI
        \end{verbatim}
        \begin{lstlisting}[language=C]
if (temperature < 19.0)
{
    printf("Chauffage ON\n");
}
else
{
    printf("Chauffage OFF\n");
}
        \end{lstlisting}
    \end{UPSTIcorrectionP}
\end{UPSTIprofOnlyEnv}


\begin{UPSTIexercice}{Traduire du C vers le pseudo-code}
    On considère le programme suivant, qui affiche un message selon le niveau d'une batterie (\texttt{niveau}, en pourcentage).
    \begin{lstlisting}[language=C]
if (niveau < 15)
{
    printf("Batterie faible\n");
}
else if (niveau < 90)
{
    printf("Batterie en charge normale\n");
}
else
{
    printf("Batterie pleine\n");
}
    \end{lstlisting}
    \UPSTIquestion{Traduire ce programme en pseudo-code, en utilisant les mots-clés \texttt{SI}, \texttt{SINON SI}, \texttt{SINON}, \texttt{FIN SI} et \texttt{AFFICHER}.}

\end{UPSTIexercice}

\begin{UPSTIprofOnlyEnv}
    \begin{UPSTIcorrectionP}{Traduire du C vers le pseudo-code}
        \begin{verbatim}
SI niveau < 15 ALORS
    AFFICHER "Batterie faible"
SINON SI niveau < 90 ALORS
    AFFICHER "Batterie en charge normale"
SINON
    AFFICHER "Batterie pleine"
FIN SI
        \end{verbatim}
    \end{UPSTIcorrectionP}
\end{UPSTIprofOnlyEnv}
